\section{Introduction}
\label{sec:introduction}

A coding agent accumulates evidence as it works. A command exposes a
dependency, an edit changes the meaning of an earlier observation, and a
test rules out a proposed fix. An interaction transcript preserves these
events together with their context. Replacing that transcript with compact
state requires deciding which information later decisions will need.
The consequences extend beyond the length of the next prompt.

SKILL.state makes this replacement an explicit runtime
design~\citep{skillstate2026}. Persistent instructions, mutable structured
state, and the current observation replace an expanding conversation.
The attraction is clear: when those inputs are bounded, context need not
grow with the trajectory. But a bound on context is not a bound on the
work required to solve the task. A smaller representation may omit evidence,
require additional state-writing output, or change how the agent acts and
stops. Cache-aware pricing further separates prompt size from monetary cost.

We study this tension in a terminal-based coding agent. The evaluation
transfers the state-based runtime idea to a different model, agent, and
benchmark; it is not a reproduction of the original study. Our central
question is what information must survive history replacement for
compression to remain useful over a complete task.

The comparisons separate decisions that are easily conflated under
agent memory. First, who writes the memory? A model can maintain a
semantic account of its work, while the runtime can retain observed tool
events directly. This motivates hybrid state: execution records and
model-written interpretation have different owners. Second, what does the
runtime retain? File references, modifications, command outcomes, and
diagnostics convey different amounts of execution evidence. Third, which
original interactions remain available? A recent-history window preserves
detail without requiring the state writer to anticipate its relevance.
Finally, the update protocol and schema determine how the model expresses
its semantic memory.

The results connect these design choices to task success and resource use.
Model-written state substantially reduces peak prompts while lowering
success and increasing conditional successful-execution cost. An
exploratory component comparison shows why hybrid memory cannot be treated
as a single intervention: richer execution records perform better than
file references alone, but additional records do not yield continuous
improvement. Retaining recent interactions produces the clearest recovery
in both state families. The tested schema and protocol changes do not.

These findings motivate evaluating structured state together with its
retained evidence and surrounding history. We distinguish partial recovery from a complete
replacement for history, and separate measured results from proposed
adaptive switching or archival retrieval. The practical outcome is a set
of quality--cost trade-offs, not a universal recommendation to discard or
retain every transcript.
